% ======================================================================
\section{Prescriptions and distributions}\label{app:distributions}
% ======================================================================

We collect here the conventions used for the singular longitudinal and transverse integrals.

\paragraph{Principal values.} Energy denominators that vanish inside the region of integration, such as $p^+\v k^2-k^+\v p^2$ in momentum space or $p^+(\v x-\v u)^2-k^+(\v x-\v z)^2$ in coordinate space, arise from transitions between states of equal free energy and are understood as Cauchy principal values. Correspondingly, logarithms that arise from integrals over such denominators are logarithms of absolute values, $\log|\cdot|$. This is the standing-wave dressing of the final state; the out-state dressing differs by the finite term $\Delta N$ of Eq.~\eqref{eq:DeltaN} (Section~\ref{sec:fin_pv}).

\paragraph{Dimensional regularisation.} Transverse loop integrals are computed in $d=2-\epsilon$ dimensions, $\int\frac{d^2p}{(2\pi)^2}\to\mu^{\epsilon}\int\frac{d^dp}{(2\pi)^d}$, in the $\MS$ scheme with scale $\mu_{\MS}$. Scaleless integrals vanish, $\int d^dp/\v p^2=0$. Coordinate-space logarithms are written with the combination~\eqref{eq:c0}.

\paragraph{Plus prescriptions in $p^+$.} For a function $f(p^+)$ regular at $p^+=0$,
\begin{align}\label{eq:plusp}
\int_0^{P}dp^+\Big[\frac{1}{p^+}\Big]_+f(p^+)\equiv\int_0^{P}dp^+\,\frac{f(p^+)-f(0)}{p^+}\,.
\end{align}
The subtracted term, integrated from $\Lambda$, is $f(0)\log(P/\Lambda)$; when $f(0)$ is the JIMWLK (target-side) structure it is part of the JIMWLK term.

\paragraph{Plus prescriptions in the momentum fraction.} With $\xi=k^+/(p^++k^+)$ the momentum fraction of the observed gluon in its parent, and for a function $F(\xi)$,
\begin{align}\label{eq:plusxi}
\int_0^1d\xi\,\frac{\xi}{(1-\xi)_+}F(\xi)&\equiv\int_0^1d\xi\,\frac{\xi F(\xi)-F(1)}{1-\xi}\,,\notag\\
\int_0^1d\xi\,\Big[\frac{1-\xi}{\xi}\Big]_+F(\xi)&\equiv\int_0^1d\xi\,\frac{1-\xi}{\xi}\big[F(\xi)-F(0)\big]\,.
\end{align}
The first is the standard prescription at $\xi\to1$ (integrated gluon soft); the second subtracts the pole at $\xi\to0$ (observed gluon soft relative to its parent), which belongs to the projectile (BFKL) evolution. The combination of the two subtractions coincides with the ``double plus'' prescription of~\cite{Kovner:2023vsy}: with $N(\xi)=\xi^2+(1-\xi)^2+\xi^2(1-\xi)^2$, $N(0)=N(1)=1$,
\begin{align}
\int_0^1\frac{d\xi}{\xi_+(1-\xi)_+}\,N(\xi)\,F(\xi)&\equiv\int_0^1d\xi\Big[\frac{N(\xi)F(\xi)}{\xi(1-\xi)}-\frac{F(0)}{\xi}-\frac{F(1)}{1-\xi}\Big]\,,\notag\\
\frac{N(\xi)}{\xi(1-\xi)}&=\xi(1-\xi)+\frac{\xi}{1-\xi}+\frac{1-\xi}{\xi}\,,
\end{align}
and $\int_0^1\frac{d\xi}{\xi_+(1-\xi)_+}N(\xi)=-\frac{11}{6}$.

\paragraph{Integrals over $p^+$ up to $\vee\to\infty$.} Some terms have no dependence on $p^+$ other than an oscillating phase. Their integral converges for $\vee\to\infty$ in the sense of distributions,
\begin{align}\label{eq:distrib}
\int_0^\infty dp^+\,e^{-iap^+}=\pi\,\delta(a)-i\,\mathcal P\frac1a\,.
\end{align}

\paragraph{Recoil function.} The recoil of the observed gluon in the rows with $\mf B_3$ and $\mf C_1$ produces, after the longitudinal integration, the function
\begin{align}\label{eq:Erecoil}
\mathcal E(s)=\int_0^1\frac{dt}{t}\big(e^{-ist}-1\big)\,,
\end{align}
which vanishes at $s=0$ and grows logarithmically for $|s|\to\infty$, $\mathcal E(s)=-\log|s|-\gamma_E-\frac{i\pi}{2}{\rm sign}(s)+O(1/s)$.

\paragraph{Longitudinal integrals.} The following integrals, with $\delta=\Lambda/k^+\to0$, produce the constants of the one-gluon wave function at order $g^3$, Eq.~\eqref{eq:xiintegrals}:
\begin{align}
\int_\delta^{1-\delta}d\xi\Big[-2+\xi(1-\xi)+\frac{1}{\xi(1-\xi)}\Big]&=-\frac{11}{6}-2\log\delta\,,\notag\\
\int_0^1d\xi\,\big[-2+\xi(1-\xi)\big]\log\big(\xi(1-\xi)\big)&=\frac{67}{18}\,,\notag\\
\int_\delta^{1-\delta}d\xi\,\frac{\log(\xi(1-\xi))}{\xi(1-\xi)}&=-\log^2\delta-\frac{\pi^2}{3}\,,
\end{align}
up to terms that vanish for $\delta\to0$.
